\begin{abstract}
How efficiently can an external controller steer a group of interacting
language-model agents? We study this question in relational reasoning tasks
where evidence is distributed across the population. Our feedback controller
has two separate limits: it observes only a sample of the agents and can
intervene in only a fixed number of their interactions. This setup lets us
measure the distinct effects of sensing and actuation. We introduce finite-time
measures for three aspects of control: how much an intervention changes the
population, how much information about the intervention appears in the next
population state, and how efficiently that information produces directed
change. We also derive a finite-state model that separates the direction of
control from the probability that an agent responds. In experiments with
matched initial social and knowledge states, advocacy increases support for the
controller's target in every tested condition, with susceptibilities from
$0.100$ to $0.163$. The transmitted control information depends strongly on
the population state. More actuation produces a wider informative region,
whereas stronger feedback gain is predicted to concentrate information near
the controller's decision boundary. Information--response efficiency increases
with actuation across the full resource grid, from about $0.4\%$ to $2.5\%$ of
the Pinsker upper bound. A calibrated finite-compliance model also reproduces
the observed controlled current closely ($r=0.999$). Finally, an exploratory
thermodynamic-efficiency measure shows diminishing returns from additional
sensing when the added information does not produce proportional population
change. In the regime studied here, actuation is therefore more useful than
additional sensing. Our framework measures control efficiency while keeping it
separate from whether the agents reach the correct answer.
\end{abstract}
